
\documentclass{article}
\usepackage{colortbl}
\usepackage{makecell}
\usepackage{multirow}
\usepackage{supertabular}

\begin{document}

\newcounter{utterance}

\centering \large Interaction Transcript for game `hot\_air\_balloon', experiment `air\_balloon\_survival\_et\_complexity\_hard', episode 0 with gpt{-}5.4{-}azure.
\vspace{24pt}

{ \footnotesize  \setcounter{utterance}{1}
\setlength{\tabcolsep}{0pt}
\begin{supertabular}{c@{$\;$}|p{.15\linewidth}@{}p{.15\linewidth}p{.15\linewidth}p{.15\linewidth}p{.15\linewidth}p{.15\linewidth}}
    \# & $\;$A & \multicolumn{4}{c}{Game Master} & $\;\:$B\\
    \hline

    \theutterance \stepcounter{utterance}  
    & & \multicolumn{4}{p{0.6\linewidth}}{
        \cellcolor[rgb]{0.9,0.9,0.9}{
            \makecell[{{p{\linewidth}}}]{
                \texttt{\tiny{[P1$\langle$GM]}}
                \texttt{Te osalete ühises läbirääkimismängus.} \\
\\ 
\texttt{Te peate koos teise osalejaga kokku leppima ühes esemete kogumis, mida hoitakse. Igaühel teist on oma seisukoht, kui oluline on iga ese teie jaoks (tähtsus). Te ei tea, kuidas teine osaleja väärtustab esemeid. Lisaks sellele on teile antud, kui suurt pingutust iga ese nõuab.} \\
\texttt{Te võite kokku leppida ainult siis, kui valitud esemete kogupingutus ei ületa ühist piirmäära:} \\
\\ 
\texttt{LIMIT = 4734} \\
\\ 
\texttt{Siin on üksikute esemete pingutuse väärtused:} \\
\\ 
\texttt{Eseme pingutus = \{"C08": 192, "C55": 325, "A43": 812, "A67": 868, "B42": 593, "A62": 915, "B19": 944, "C43": 311, "C73": 252, "C10": 343, "B36": 104, "A18": 558, "C61": 627, "C67": 593, "C42": 827, "C70": 611, "B40": 95, "C39": 251, "A21": 226, "C16": 21\}} \\
\\ 
\texttt{Siin on teie isiklik seisukoht iga eseme tähtsuse kohta:} \\
\\ 
\texttt{Eseme tähtsuse väärtused = \{"C08": 286, "C55": 419, "A43": 906, "A67": 962, "B42": 687, "A62": 1009, "B19": 1038, "C43": 405, "C73": 346, "C10": 437, "B36": 198, "A18": 652, "C61": 721, "C67": 687, "C42": 921, "C70": 705, "B40": 189, "C39": 345, "A21": 320, "C16": 115\}} \\
\\ 
\texttt{Eesmärk:} \\
\\ 
\texttt{Teie eesmärk on pidada läbirääkimisi ühise esemete kogumi üle, mis toob teile võimalikult palju kasu (st maksimeerib kogutähtsust TEIE jaoks), jäädes samal ajal LIMITi piiridesse. Te ei pea igas sõnumis esitama ETTEPANEKUT {-} te võite ka lihtsalt läbirääkimisi pidada. Kõik taktikad on lubatud!} \\
\\ 
\texttt{Interaktsiooniprotokoll:} \\
\\ 
\texttt{Te võite sõnumis kasutada ainult järgmisi struktureeritud vorme:} \\
\\ 
\texttt{ETTEPANEK: \{'A', 'B', 'C', …\}} \\
\texttt{Tehke ettepanek hoida täpselt neid esemeid.} \\
\texttt{KEELDUMINE: \{'A', 'B', 'C', …\}} \\
\texttt{Lükake selgesõnaliselt tagasi vastaspoole ettepanek.} \\
\texttt{ARGUMENT: \{'...'\}} \\
\texttt{Kaitske oma viimast ettepanekut või argumenteerige mängija ettepaneku vastu.} \\
\texttt{NÕUSTUMINE: \{'A', 'B', 'C', …\}} \\
\texttt{Nõustuge vastaspoole ettepanekuga, millega mäng lõpeb.} \\
\\ 
\\ 
\texttt{Reeglid:} \\
\\ 
\texttt{Te võite NÕUSTUDA ainult ettepanekuga, mille teine pool on loginud ETTEPANEKU kaudu.} \\
\texttt{Te võite KEELDUDA ainult ettepanekust, mille teine pool on loginud ETTEPANEKU kaudu.} \\
\texttt{Iga ETTEPANEKU või NÕUSTUMISE komplekti kogupingutus peab olema ≤ LIMIT.} \\
\texttt{ÄRGE avaldage oma varjatud tähtsuse skoore.} \\
\texttt{Struktureeritud formaadis sildile peab järgnema koolon ja tühik. Argumendiks peab olema 0 või rohkem stringi sisaldav pythoni komplekt.} \\
\texttt{Seega peab see olema kujul "TAG: \{...\}"} \\
\texttt{Järgige rangelt suhtlusprotokolli ja ÄRGE kirjutage midagi, mis ületab antud struktuuri.} \\
\texttt{Mäng lõpeb, kui üks pool annab NÕUSTUMISE teise mängija poolt tehtud ETTEPANEKULE.} \\
\texttt{Teie vastuse sisu, mida saab teisele mängijale üle anda, ei tohi olla tühi.} \\
\texttt{Aktiivsed on ainult need ettepanekud, mis on kirja pandud ETTEPANEKU vormingus struktuuri kaudu ja millest ei ole KEELDUMISE kaudu keeldutud.} \\
\\ 
            }
        }
    }
    & \\ \\

    \theutterance \stepcounter{utterance}  
    & \multicolumn{4}{p{0.6\linewidth}}{
        \cellcolor[rgb]{0.9,0.9,0.9}{
            \makecell[{{p{\linewidth}}}]{
                \texttt{\tiny{[P1$\rangle$GM]}}
                \texttt{ETTEPANEK: \{'C08', 'C55', 'A67', 'B42', 'C43', 'C73', 'C10', 'B36', 'A18', 'C67', 'B40', 'C39', 'A21', 'C16'\}} \\
            }
        }
    }
    & & \\ \\

    \theutterance \stepcounter{utterance}  
    & & \multicolumn{4}{p{0.6\linewidth}}{
        \cellcolor[rgb]{0.9,0.9,0.9}{
            \makecell[{{p{\linewidth}}}]{
                \texttt{\tiny{[GM$\rangle$P2]}}
                \texttt{Te osalete ühises läbirääkimismängus.} \\
\\ 
\texttt{Te peate koos teise osalejaga kokku leppima ühes esemete kogumis, mida hoitakse. Igaühel teist on oma seisukoht, kui oluline on iga ese teie jaoks (tähtsus). Te ei tea, kuidas teine osaleja väärtustab esemeid. Lisaks sellele on teile antud, kui suurt pingutust iga ese nõuab.} \\
\texttt{Te võite kokku leppida ainult siis, kui valitud esemete kogupingutus ei ületa ühist piirmäära:} \\
\\ 
\texttt{LIMIT = 4734} \\
\\ 
\texttt{Siin on üksikute esemete pingutuse väärtused:} \\
\\ 
\texttt{Eseme pingutus = \{"C08": 192, "C55": 325, "A43": 812, "A67": 868, "B42": 593, "A62": 915, "B19": 944, "C43": 311, "C73": 252, "C10": 343, "B36": 104, "A18": 558, "C61": 627, "C67": 593, "C42": 827, "C70": 611, "B40": 95, "C39": 251, "A21": 226, "C16": 21\}} \\
\\ 
\texttt{Siin on teie isiklik seisukoht iga eseme tähtsuse kohta:} \\
\\ 
\texttt{Eseme tähtsuse väärtused = \{"C08": 286, "C55": 419, "A43": 906, "A67": 962, "B42": 687, "A62": 1009, "B19": 1038, "C43": 405, "C73": 346, "C10": 437, "B36": 198, "A18": 652, "C61": 721, "C67": 687, "C42": 921, "C70": 705, "B40": 189, "C39": 345, "A21": 320, "C16": 115\}} \\
\\ 
\texttt{Eesmärk:} \\
\\ 
\texttt{Teie eesmärk on pidada läbirääkimisi ühise esemete kogumi üle, mis toob teile võimalikult palju kasu (st maksimeerib kogutähtsust TEIE jaoks), jäädes samal ajal LIMITi piiridesse. Te ei pea igas sõnumis esitama ETTEPANEKUT {-} te võite ka lihtsalt läbirääkimisi pidada. Kõik taktikad on lubatud!} \\
\\ 
\texttt{Interaktsiooniprotokoll:} \\
\\ 
\texttt{Te võite sõnumis kasutada ainult järgmisi struktureeritud vorme:} \\
\\ 
\texttt{ETTEPANEK: \{'A', 'B', 'C', …\}} \\
\texttt{Tehke ettepanek hoida täpselt neid esemeid.} \\
\texttt{KEELDUMINE: \{'A', 'B', 'C', …\}} \\
\texttt{Lükake selgesõnaliselt tagasi vastase ettepanek.} \\
\texttt{ARGUMENT: \{'...'\}} \\
\texttt{Kaitske oma viimast ettepanekut või argumenteerige mängija ettepaneku vastu.} \\
\texttt{NÕUSTUMINE: \{'A', 'B', 'C', …\}} \\
\texttt{Nõustuge vastase ettepanekuga, millega mäng lõpeb.} \\
\\ 
\\ 
\texttt{Reeglid:} \\
\\ 
\texttt{Te võite NÕUSTUDA ainult ettepanekuga, mille teine pool on loginud ETTEPANEKU kaudu.} \\
\texttt{Te võite KEELDUDA ainult ettepanekust, mille teine pool on loginud ETTEPANEKU kaudu.} \\
\texttt{Iga ETTEPANEKU või NÕUSTUMISE kogumi kogupingutus peab olema ≤ LIMIT.} \\
\texttt{ÄRGE avaldage oma varjatud tähtsuse skoore.} \\
\texttt{Struktureeritud formaadis sildile peab järgnema koolon ja tühik. Argumendiks peab olema 0 või rohkem stringi sisaldav pythoni kogum.} \\
\texttt{Seega peab see olema kujul "TAG: \{...\}"} \\
\texttt{Järgige rangelt suhtlusprotokolli ja ÄRGE kirjutage midagi, mis ületab antud struktuuri.} \\
\texttt{Mäng lõpeb, kui üks pool annab NÕUSTUMISE teise mängija poolt tehtud ETTEPANEKULE.} \\
\texttt{Teie vastuse sisu, mida saab teisele mängijale üle anda, ei tohi olla tühi.} \\
\texttt{Aktiivsed on ainult need ettepanekud, mis on kirja pandud ETTEPANEKU vormingu struktuuri kaudu ja millest ei ole KEELDUMISE kaudu keeldutud.} \\
\\ 
\\ 
\\ 
\texttt{Saate nüüd teise mängija esimese sõnumi.} \\
\\ 
\texttt{ETTEPANEK: \{'C08', 'C55', 'A67', 'B42', 'C43', 'C73', 'C10', 'B36', 'A18', 'C67', 'B40', 'C39', 'A21', 'C16'\}} \\
            }
        }
    }
    & \\ \\

    \theutterance \stepcounter{utterance}  
    & & & \multicolumn{4}{p{0.6\linewidth}}{
        \cellcolor[rgb]{0.9,0.9,0.9}{
            \makecell[{{p{\linewidth}}}]{
                \texttt{\tiny{[GM$\langle$P2]}}
                \texttt{NÕUSTUMINE: \{'C08', 'C55', 'A67', 'B42', 'C43', 'C73', 'C10', 'B36', 'A18', 'C67', 'B40', 'C39', 'A21', 'C16'\}} \\
            }
        }
    }
     \\ \\

    \theutterance \stepcounter{utterance}  
    & & & \multicolumn{2}{p{0.3\linewidth}}{
        \cellcolor[rgb]{0.9,0.9,0.9}{
            \makecell[{{p{\linewidth}}}]{
                \texttt{\tiny{[GM$|$GM]}}
                \texttt{game successful} \\
            }
        }
    }
    & & \\ \\

    \theutterance \stepcounter{utterance}  
    & & & \multicolumn{2}{p{0.3\linewidth}}{
        \cellcolor[rgb]{0.9,0.9,0.9}{
            \makecell[{{p{\linewidth}}}]{
                \texttt{\tiny{[GM$|$GM]}}
                \texttt{end game} \\
            }
        }
    }
    & & \\ \\

\end{supertabular}
}

\end{document}
